\documentclass[10pt]{article}
\usepackage[T1]{fontenc}
\usepackage[utf8]{inputenc}
\usepackage[french]{babel}
\usepackage{amsmath,amssymb,mathrsfs,mathpazo}
\usepackage{geometry}
\geometry{a4paper,margin=2cm}
\usepackage{graphicx,array,tabularx,tabularray,colortbl,multirow,multicol,booktabs,enumitem,pifont,ulem,fancybox,diagbox,cancel}
\usepackage{pstricks,pstricks-add,pst-plot,pst-tree,pst-node,pst-text,pst-func,pst-math,pst-bspline,pst-3dplot,pst-eucl,pst-solides3d}
\let\pscalculate\undefined
\usepackage{pst-bezier}
\usepackage{tikz,pgfplots,tkz-tab}
\usetikzlibrary{arrows,arrows.meta,patterns,calc,positioning,trees,decorations.markings,shapes.geometric,shapes.misc,shapes.arrows,intersections,decorations.pathmorphing,decorations.pathreplacing,backgrounds,shadows,babel}
\pgfplotsset{compat=1.18}
\usepgfplotslibrary{fillbetween}
\usepackage[np]{numprint}
\usepackage{listings,makecell,pas-tableur,twoopt}
\let\barre\undefined
\usepackage[pstricks]{tablvar}
\definecolor{darkgreen}{RGB}{0,100,0}
\providecommand{\red}{\color{red}}
\providecommand{\blue}{\color{blue}}
\providecommand{\green}{\color{green}}
\providecommand{\black}{\color{black}}
\providecommand{\white}{\color{white}}
\providecommand{\gray}{\color{gray}}
\providecommand{\cyan}{\color{cyan}}
\providecommand{\magenta}{\color{magenta}}
\providecommand{\yellow}{\color{yellow}}
\providecommand{\violet}{\color{violet}}
\providecommand{\orange}{\color{orange}}
\providecommand{\euro}{\texteuro}
\providecommand{\EUR}{\texteuro}
\providecommand{\ieme}{\textsuperscript{e}}
\providecommand{\degres}{\ensuremath{^\circ}}
\providecommand{\degree}{\ensuremath{^\circ}}
\providecommand{\vv}[1]{\overrightarrow{#1}}
\providecommand{\Coord}[1]{#1}
\providecommand{\decosix}{}
\providecommand{\starredbullet}{\ensuremath{\bullet}}
\pagestyle{empty}
\setlength{\parindent}{0pt}
\usepackage{adjustbox}\DeclareUnicodeCharacter{2236}{:}\DeclareUnicodeCharacter{2190}{\ensuremath{\leftarrow}}
\providecommand{\R}{\mathbb{R}}
\providecommand{\N}{\mathbb{N}}
\providecommand{\Z}{\mathbb{Z}}
\providecommand{\Q}{\mathbb{Q}}
\providecommand{\D}{\mathbb{D}}
\providecommand{\C}{\mathbb{C}}
\providecommand{\e}{{\mathrm{e}}}
\providecommand{\vect}[1]{\overrightarrow{#1}}
\providecommand{\vectt}[1]{\overrightarrow{\mathrm{#1}}}
\providecommand{\vvt}[1]{\overrightarrow{\mathrm{#1}}}
\providecommand{\np}[2][]{{#2~#1}}
\providecommand{\vv}[1]{{\overrightarrow{#1}}}
\renewcommand{\d}{{\mathrm{d}}}
\providecommand{\degres}{{^\circ}}
\providecommand{\degre}{{^\circ}}
\providecommand{\blue}{\color{blue}}
\providecommand{\red}{\color{red}}
\providecommand{\enskip}{\quad}
\providecommand{\llbracket}{\lbrack\!\lbrack}
\providecommand{\rrbracket}{\rbrack\!\rbrack}
\providecommand{\euro}{{\text{€}}}
\providecommand{\no}{{\text{n°}}}
\providecommand{\interval}[2]{{[#1;#2]}}
\providecommand{\intervFO}[2]{{[#1;#2[}}
\providecommand{\dotfill}{\ldots}

\providecommand{\ds}{\displaystyle}
\providecommand{\textding}[1]{#1}
\providecommand{\Oij}{\(\mathrm{O};\vec{\imath},\vec{\jmath}\)}
\providecommand{\Oijk}{\(\mathrm{O};\vec{\imath},\vec{\jmath},\vec{k}\)}
\providecommand{\Ouv}{\(\mathrm{O};\vec{u},\vec{v}\)}
\providecommand{\intervFF}[2]{\ensuremath{[#1;#2]}}
\renewcommand{\euro}{\ifmmode\text{\texteuro}\else\texteuro\fi}
\DeclareUnicodeCharacter{00B0}{\ifmmode^{\circ}\else\textdegree\fi}
\DeclareUnicodeCharacter{20AC}{\ifmmode\text{\texteuro}\else\texteuro\fi}
% Énoncés : Éducation nationale. Transcription et corrigés : APMEP, Denis Vergès et contributeurs. Découpage et classement : CoopMaths, Sébastien Lozano, Éric Elter et contributeurs. Adaptation de la présentation : Maths Explorer.
% CC BY-SA : https://coopmaths.fr/www/about/
% https://coopmaths.fr/alea/static/bac/2021/tex/bac_2021_05_ameriquenord_2_cor.tex
\begin{document}
\noindent\textit{APMEP / CoopMaths. CC BY-SA. Présentation adaptée par Maths Explorer.}\par\medskip
\medskip









\[\left\{\begin{array}{l c l}
u_0		&=	&0,6\\
u_{n+1}	&=	&0,75 u_n\left(1 - 0,15 u_n\right)
\end{array}\right.\]

où pour tout entier naturel $n$, $u_n$ désigne le nombre d'individus, en milliers, au début de l'année $2020 + n$.

\medskip

\begin{enumerate}
\item ~
$\bullet~~$2021 correspond à $n = 1$, donc $u_1 = 0,75u_0 \times \left(1 - 0,15u_0\right) = 0,75 \times 0,6\times (1 - 0,15 \times 0,6) = 0,45 \times (1 - 0,09) = 0,45 \times0,91 = \np{0,4095}$ soit environ $410$ individus.

$\bullet~~$2022 correspond à $n = 2$, donc $u_2 = 0,75u_1 \times \left(1 - 0,15u_1\right) = 0,75 \times \np{0,4095}\times (1 - 0,15 \times \np{0,4095}) = \np{0,307125} \times (1 - \np{0,061425}) = \np{0,307125} \times \np{0,938575}) \approx \np{0,2282}$ soit environ 228~individus.

\end{enumerate}

Soit $f$ la fonction définie sur l'intervalle [0~;~1] par 

\[f(x) = 0,75x (1 - 0,15x).\]

\begin{enumerate}[resume]
\item 
$f$ est une fonction polynôme dérivable sur $\R$, donc sur [0~;~1] et sur cet intervalle :

$f'(x)  = 0,75(1 - 0,15x) - 0,75x \times 0,15 = 0,75 - \np{0,1125}x - \np{0,1125}x = 0,75 - 0,225x$.

Or $0 \leqslant x \leqslant 1 \Rightarrow 0 \leqslant 0,225x \leqslant 0,225 \Rightarrow - 0,225 \leqslant -0,225x \leqslant 0 \Rightarrow $

$0,75 - 0,225 \leqslant 0,75 - 0,225x \leqslant 0,75$ ou enfin $0,525 \leqslant f'(x) \leqslant 0,75$.

Sur [0~;~1], $f'(x) > 0$, donc $f$est strictement croissante de $f(0) = 0$ à $f(1) = 0,75 \times 0,85 = \np{0,6375}$.
\item 

Sur [0~;~1], \, $f(x) = x \iff 0,75x (1 - 0,15x) = x \iff 0,75x(1 - 0,15x) - x = 0 \iff x[0,75(1 - 0,15x) - 1] = 0 \iff x(0,75 - \np{0,1125}x - 1) = 0 \iff x(- 0,25 - \np{0,1125}x) = 0 \iff \left\{\begin{array}{l c l}
x&=&0\,\text{ou}\\
-0,25 - \np{0,1125}x&=0&
\end{array}\right. \iff \left\{\begin{array}{l c l}
x&=&0\,\text{ou}\\
-0,25 &=&\np{0,1125}x
\end{array}\right. \iff \left\{\begin{array}{l c l}
x&=&0\,\text{ou}\\
-\frac{0,25}{{0,1125}}&=&x
\end{array}\right. $.

Or $-\frac{0,25}{{0,1125}} < 0$ donc dans [0~;~1], \, $S = \{0\}$.
\end{enumerate}

On remarquera pour la suite de l'exercice que, pour tout entier naturel $n$,\, $u_{n+1} = f\left(u_n\right)$.

\begin{enumerate}[resume]
\item 
	\begin{enumerate}
		\item ~
\emph{Initialisation} : on a vu que $0 \leqslant \np{0,4095} \leqslant 0,6 \leqslant 1$, soit $0 \leqslant u_1 \leqslant u_0\leqslant 1$ : la relation est vraie au rang $0$ ;

\emph{Hérédité} : Supposons que pour $n \in \N$, on ait :

$0 \leqslant u_{n+1} \leqslant  u_n \leqslant 1$ ; la fonction $f$ étant strictement croissante sur [0~;~1], on a donc :
$f(0) \leqslant f\left(u_{n+1}\right) \leqslant f\left(u_n\right) \leqslant f(1)$, 

soit puisque $f(0) = 0$ et $f(1) = 0,75 \times (1 - 0,15) = \np{0,6375} \leqslant 1$ :

$0 \leqslant u_{n+2} \leqslant u_{n+1} \leqslant  1$ : la relation est donc vraie au rang $n + 1$.

\textbf{Conclusion :} la relation est vraie au rang 0 et si elle est vraie au rang $n$ naturel quelconque, elle est vraie au rang $n+1$ : d'après le principe de récurrence :

Pour tout entier naturel $n$,\, $0 \leqslant u_{n+1} \leqslant  u_n \leqslant 1$.	
		
		\item 
La suite $\left(u_n\right)$ est d'après la question précédente décroissante et minorée par $0$ ; elle est donc est convergente.
		\item 
Le résultat précédent montre que la suite $\left(u_n\right)$ converge vers un nombre $\ell \geqslant 0$ et ce nombre $\ell$ vérifie l'équation $f(x) = x$, dont on a vu à la question \textbf{3.} qu'elle n'avait que $0$ comme solution.
		
Conclusion : $\displaystyle\lim_{n \to + \infty} u_n = \ell = 0$.
	\end{enumerate}
\item 
	\begin{enumerate}
		\item 
L'étude précédente a montré que le nombre d'individus décroit, donc le biologiste a raison puisque la limite de la suite du nombre d'individus est égale à zéro.
		\item 
















L'algorithme calcule les termes de la suite tant que ceux-ci sont supérieurs à 0,02

Il s'arrête à $n = 11$ car $u_{10} \approx 0,019$

L'espèce sera donc menacée d'extinction en 2031.

	\end{enumerate}
\end{enumerate}

\end{document}
